\chapter{Loops}\label{ch:loops}

\begin{goals}
\begin{itemize}
  \item Repeating a task a fixed number of times with \code{repeat}
  \item Counting from one number to another, upwards, downwards and in steps
  \item Repeating while a condition holds, with \code{until}
  \item Breaking out of a loop with \code{break} and skipping a round with \code{continue}
  \item A loop inside a loop
\end{itemize}
\end{goals}

\section{Why loops?}

Suppose you want to print the numbers 1 to 5. With what you have learned so
far, you would write five \code{println} instructions. Now what if you want
to print up to 1000? Or if the count is not known in advance? Computers were
built precisely for this kind of repetitive work, and a \textbf{loop} is the
tool that tells them: ``do this many times over.''

\section{Repeating a fixed number of times}

The simplest loop in Salam is \kwd{repeat}. You write after it how many
times, and Salam runs the block that many times:

\program{The simplest loop}
\begin{salamcode}
func main:
    repeat 3:
        println "Hello!"
    end
    println "Done"
end
\end{salamcode}

\begin{salamout}
Hello!
Hello!
Hello!
Done
\end{salamout}

The structure of a loop is the familiar one: a line ending in a colon, an
indented block, and \code{end}. Each run of the block is called an
\textbf{iteration}. This loop has three iterations. The number after
\code{repeat} can also be a variable or a calculation:
\code{repeat guests:} or \code{repeat n * 2:}.

\section{The loop counter}

Most of the time we want to know which iteration we are in. By writing
\kwd{in} and a name, you ask Salam to put the number of the iteration into a
variable each time round. The numbering starts at \emph{zero}:

\program{The loop counter}
\begin{salamcode}
func main:
    repeat 5 in i:
        println "Iteration", i
    end
end
\end{salamcode}

\begin{salamout}
Iteration 0
Iteration 1
Iteration 2
Iteration 3
Iteration 4
\end{salamout}

Five iterations, numbered 0 to 4. Starting at zero may seem odd, but it is a
very old custom in programming, and in Chapter~\ref{ch:arrays} you will see
why it is useful. The variable \code{i} exists only inside the loop; after
the \code{end} it is no longer available.

\section{Counting from one number to another}

If you want to count from 1 rather than 0, or from 10 to 20, use the form
\code{repeat start to finish}. Both ends of the range are counted:

\program{Counting from one to five}
\begin{salamcode}
func main:
    repeat 1 to 5 in i:
        println i, "times 3 is", i * 3
    end
end
\end{salamcode}

\begin{salamout}
1 times 3 is 3
2 times 3 is 6
3 times 3 is 9
4 times 3 is 12
5 times 3 is 15
\end{salamout}

If the starting number is larger than the finishing one, Salam counts
downwards:

\program{A countdown}
\begin{salamcode}
func main:
    repeat 5 to 1 in second:
        println second
    end
    println "Lift-off!"
end
\end{salamcode}

\begin{salamout}
5
4
3
2
1
Lift-off!
\end{salamout}

And with \kwd{by} you can count several at a time instead of one by one:

\program{Counting in steps}
\begin{salamcode}
func main:
    mut line := ""
    repeat 0 to 20 by 5 in n:
        line = line + n + " "
    end
    println line
end
\end{salamcode}

\begin{salamout}
0 5 10 15 20
\end{salamout}

In this program, instead of printing each number on its own line, we lined
them all up in one string and printed it once at the end. You will see this
trick (start with an empty string and attach the pieces) many times. The
step must be a positive number; to count down in steps, simply make the
start larger than the finish.

\section{Adding up in a loop}

Now let us combine the ``adding up'' pattern of Chapter~\ref{ch:variables}
with a loop. We make a changeable variable with the value zero and add
something to it in each iteration:

\program{The sum of one to a hundred}
\begin{salamcode}
func main:
    mut total := 0
    repeat 1 to 100 in n:
        total += n
    end
    println "Sum of 1 to 100:", total
end
\end{salamcode}

\begin{salamout}
Sum of 1 to 100: 5050
\end{salamout}

Notice that \code{total} must be \code{mut}, since it changes in every
iteration, and it must be declared \emph{outside} the loop; had we declared
it inside, a fresh box holding zero would have been made in every iteration
and nothing would ever have added up.

Loops and conditions together make a very powerful combination. For
example, adding up and counting only the even numbers:

\program{Counting and adding the even numbers}
\begin{salamcode}
func main:
    mut even total := 0
    mut even count := 0
    repeat 1 to 20 in n:
        if (n % 2) == 0:
            even total += n
            even count++
        end
    end
    println "Count:", even count, "Sum:", even total
end
\end{salamcode}

\begin{salamout}
Count: 10 Sum: 110
\end{salamout}

\section{Repeating while a condition holds: until}

All the loops above knew in advance how many iterations they had. But
sometimes we do not know: ``keep saving until you reach a thousand.'' ``keep
dividing the number by ten until it becomes zero.'' For these loops there is
the word \kwd{until}. A condition follows it, and as long as the condition
is true, the block is repeated:

\program{Saving up to a goal}
\begin{salamcode}
func main:
    mut savings := 100
    mut months := 0
    until savings < 1000:
        savings += 150
        months++
    end
    println "After", months, "months the savings reached", savings
end
\end{salamcode}

\begin{salamout}
After 6 months the savings reached 1000
\end{salamout}

Salam tests the condition before \emph{every} iteration. As long as the
savings are below a thousand, it adds 150 and counts one more month. The
moment the condition becomes false, the loop is over and the program carries
on after the \code{end}. If the condition is false from the very start, the
block does not run even once.

\begin{note}
The word \code{until} may read a little unexpectedly at first: the loop
runs \emph{while} the condition is true, and stops when it becomes false.
Read \code{until savings < 1000:} as ``keep going for as long as the
savings are under a thousand''.
\end{note}

\program{Counting the digits of a number}
\begin{salamcode}
func main:
    number := 48213
    mut remaining := number
    mut digit count := 0
    until remaining > 0:
        remaining /= 10
        digit count++
    end
    println number, "has", digit count, "digits"
end
\end{salamcode}

\begin{salamout}
48213 has 5 digits
\end{salamout}

Each integer division by 10 drops one digit from the right-hand end of the
number: 48213, 4821, 482, 48, 4, 0. It took five steps to reach zero, so the
number has five digits. We made the variable \code{remaining} so that
\code{number} itself stays untouched and can be printed at the end.

\begin{warn}[The endless loop]
In an \code{until} loop, something inside the block must eventually make
the condition false; otherwise the loop runs for ever. This situation is
called an \textbf{infinite loop}, and it is one of the most common mistakes
beginners make. Salam helps you twice along the way. If you delete the line
\code{remaining /= 10} from the program above, nothing changes
\code{remaining} any more, and before running the program Salam reports
error \code{E069}: ``variable \code{'remaining'} is declared
\code{'mut'} but never mutated''. But if you mistakenly write
\code{remaining += 1} instead, the program runs and \code{remaining} never
reaches zero. The online editor stops such a program after three seconds and
writes: ``runtime error: execution timed out (possible infinite loop)''.
When that happens, look for the line that should have changed the condition.
\end{warn}

\begin{note}[Which loop?]
If the number of iterations is known in advance (or follows from a number),
\code{repeat} is the natural choice. If the number of iterations depends on
something that happens during the loop, choose \code{until}. Anything can
be done with either, but one of the two is always easier to read.
\end{note}

\section{Breaking out of a loop and skipping an iteration}

Sometimes, in the middle of a loop, we realise that there is no need to go
on. The instruction \kwd{break} ends the loop on the spot:

\program{Finding something and stopping}
\begin{salamcode}
func main:
    repeat 1 to 10 in n:
        if (n * n) > 30:
            println "The first number whose square exceeds 30:", n
            break
        end
    end
end
\end{salamcode}

\begin{salamout}
The first number whose square exceeds 30: 6
\end{salamout}

Without \code{break}, the loop would have gone on to 10 and printed the
message for 7, 8, 9 and 10 as well. \code{break} said: ``found it, that is
enough.''

The instruction \kwd{continue} is gentler: it skips only \emph{this}
iteration and moves on to the next one:

\program{Skipping the multiples of three}
\begin{salamcode}
func main:
    repeat 1 to 10 in n:
        if (n % 3) == 0: continue end
        print n, ""
    end
    println ""
end
\end{salamcode}

\begin{salamout}
1 2 4 5 7 8 10
\end{salamout}

The numbers 3, 6 and 9 were not printed, because for them \code{continue}
ran and \code{print} never got its turn. (A small trick: \code{print n, ""}
prints the number followed by a space, because a space is put between items
separated by commas. The \code{println ""} at the end closes the line.)

\section{A loop inside a loop}

Like a condition, a loop can sit inside another loop. The inner loop runs
\emph{in full} for every iteration of the outer loop. The classic example is
the multiplication table:

\program{A multiplication table}
\begin{salamcode}
func main:
    repeat 1 to 5 in row:
        mut line := ""
        repeat 1 to 5 in column:
            line = line + (row * column) + " "
        end
        println line
    end
end
\end{salamcode}

\begin{salamout}
1 2 3 4 5
2 4 6 8 10
3 6 9 12 15
4 8 12 16 20
5 10 15 20 25
\end{salamout}

The outer loop has five iterations; in each of them the inner loop goes
round five times and builds one line of the table. Twenty-five
multiplications are done in all. We declared the variable \code{line}
\emph{inside} the outer loop so that it is fresh and empty for every row.

\program{A triangle of stars}
\begin{salamcode}
func main:
    repeat 1 to 4 in row:
        mut line := ""
        repeat row:
            line += "*"
        end
        println line
    end
end
\end{salamcode}

\begin{salamout}
*
**
***
****
\end{salamout}

Here the number of iterations of the inner loop depends on the counter of
the outer loop: one star in the first row, two in the second, and so on. With
small changes to this program you can draw all sorts of shapes; the exercises
at the end of the chapter suggest a few.

\section{A more complete program}

A table converting temperatures from Celsius to Fahrenheit, from zero to
forty degrees in steps of ten:

\program{A temperature conversion table}
\begin{salamcode}
func main:
    println "Celsius -> Fahrenheit"
    println "---------------------"
    repeat 0 to 40 by 10 in celsius:
        temperature := celsius as f64
        fahrenheit := temperature * 9.0 / 5.0 + 32.0
        println celsius, "->", fahrenheit
    end
end
\end{salamcode}

\begin{salamout}
Celsius -> Fahrenheit
---------------------
0 -> 32
10 -> 50
20 -> 68
30 -> 86
40 -> 104
\end{salamout}

The loop counter is an integer, so before the floating-point calculation we
converted it with \code{as f64}. The variable \code{fahrenheit} is declared
inside the loop and without \code{mut}; in every iteration a fresh box is
made, receives a value, is printed, and disappears when the iteration ends.

\begin{summary}
\begin{itemize}
  \item \code{repeat n:} runs the block n times. With
        \code{repeat n in i:} you get the number of the iteration (from
        zero).
  \item \code{repeat start to finish in i:} counts from start to finish,
        both included; if start is larger, it counts down. \code{by} sets
        the size of the step.
  \item \code{until condition:} repeats while the condition is true.
        Something inside the loop must eventually make the condition
        false.
  \item \code{break} ends the loop; \code{continue} skips to the next
        iteration.
  \item Loops can be nested; the inner loop runs in full for every
        iteration of the outer one.
  \item A variable that accumulates in a loop must be \code{mut} and must
        be declared outside the loop.
\end{itemize}
\end{summary}

\section*{Exercises}
\markright{Exercises}

\exercise Print, on one line, the numbers from 1 to 50 that are divisible by
7.

\exercise Work out and print the product of the numbers from 1 to 10 with a
loop.

\exercise You put a sum of money in a bank at 20 percent interest a year.
With \code{until}, work out after how many years the sum doubles, printing
the balance for each year.

\exercise Change the triangle-of-stars program so that it draws the
triangle upside down (four stars first, then three, \ldots). Then write a
version that draws a hollow 5 by 5 square of stars.

\exercise With two nested loops, find and print every pair of numbers from 1
to 10 whose sum is 10 (such as 1 and 9, 2 and 8, and 5 and 5). Print each
pair only once.

\exercise How many lines does the program below print? Work it out first,
then run it:
\begin{salamcode}
func main:
    repeat 3 in a:
        repeat 4 in b:
            if b == 2: break end
            println a, b
        end
    end
end
\end{salamcode}
